\documentclass{../../kalkulus-ppt}

\author[Tetew]{Teosofi Hidayah Agung}
\date{\today}
\title[Kalkulus 1 - Bab 6]{Integrasi}
\institute[Matematika ITS]{Departemen Matematika\\ Institut Teknologi Sepuluh Nopember}
\titlegraphic{{\includegraphics[scale=0.3]{ITS.png}$\quad$\includegraphics[scale=0.02]{M.png}}}

\newcommand{\dom}{\mathcal{D}}
\newcommand{\rng}{\mathcal{R}}

\renewcommand{\arraystretch}{1.5}

\begin{document}
\begin{frame}
    \titlepage
\end{frame}

\AtBeginSection{
    {
        \begin{frame}{Daftar isi}
            \tableofcontents[currentsection]
        \end{frame}}
}

\section{Anti-turunan; Integral Taktentu}
\begin{frame}
    \frametitle{\insertsection}
    \begin{definisi}[Anti-turunan]
        Suatu fungsi $F$ disebut \textbf{anti-turunan} dari fungsi $f$ pada suatu interval $I$ jika 
        \[F'(x) = f(x)\]
        untuk semua $x$ dalam $I$.
    \end{definisi}
    \onslide<2->{
        \begin{alertblock}{Integral Taktentu}
            Himpunan semua anti-turunan dari $f$ dinotasikan dengan \textbf{integral taktentu}:
            \[\int f(x) \,dx = F(x) + C\]
            dengan $F$ adalah salah satu anti-turunan $f$ dan $C$ adalah konstanta sebarang.
        \end{alertblock}
    }
\end{frame}

\section{Integrasi dengan Substitusi}
\begin{frame}
    \frametitle{\insertsection}
    \begin{teorema}[Aturan Substitusi]
        Misalkan $u = g(x)$ adalah fungsi yang dapat diturunkan dan daerah hasilnya merupakan suatu interval $I$. Misalkan $f$ kontinu pada $I$. Maka:
        \[\int f(g(x))g'(x) \,dx = \int f(u) \,du\]
    \end{teorema}
\end{frame}

\section{Luas Sebagai Limit}
\begin{frame}
    \frametitle{\insertsection}
    \begin{definisi}[Luas Daerah]
        Luas $A$ dari daerah $S$ yang terletak di bawah grafik fungsi kontinu $f$ yang tak-negatif pada interval $[a,b]$ adalah limit dari jumlah luas persegi panjang penghampir:
        \[A = \lim_{n \to \infty} \sum_{i=1}^n f(x_i) \Delta x\]
    \end{definisi}
\end{frame}

\section{Integral Tertentu}
\begin{frame}
    \frametitle{\insertsection}
    \begin{definisi}[Integral Tertentu]
        Jika $f$ didefinisikan pada selang tertutup $[a,b]$, maka \textbf{integral tertentu} dari $f$ dari $a$ ke $b$ adalah
        \[\int_a^b f(x) \,dx = \lim_{n \to \infty} \sum_{i=1}^n f(\bar{x}_i) \Delta x_i\]
        asalkan limit ini ada. Jika limitnya ada, kita katakan bahwa $f$ \textbf{terintegralkan} pada $[a,b]$.
    \end{definisi}
\end{frame}

\section{Teorema Fundamental Kalkulus Pertama}
\begin{frame}
    \frametitle{\insertsection}
    \begin{teorema}[TFK Pertama]
        Jika $f$ kontinu pada interval tertutup $[a,b]$ dan $x$ adalah sebarang titik pada $(a,b)$, maka fungsi $F$ yang didefinisikan oleh
        \[F(x) = \int_a^x f(t) \,dt\]
        dapat diturunkan pada $(a,b)$ dan $F'(x) = f(x)$. Secara ringkas:
        \[\frac{d}{dx} \left[ \int_a^x f(t) \,dt \right] = f(x)\]
    \end{teorema}
\end{frame}

\section{Integral Tertentu dengan Substitusi}
\begin{frame}
    \frametitle{\insertsection}
    \begin{teorema}[Substitusi pada Integral Tertentu]
        Jika $g'$ kontinu pada $[a,b]$ dan $f$ kontinu pada daerah hasil $g$, maka:
        \[\int_a^b f(g(x))g'(x) \,dx = \int_{g(a)}^{g(b)} f(u) \,du\]
    \end{teorema}
\end{frame}

\section{Teorema Fundamental Kalkulus Kedua}
\begin{frame}
    \frametitle{\insertsection}
    \begin{teorema}[TFK Kedua]
        Jika $f$ kontinu pada interval tertutup $[a,b]$ dan $F$ adalah sebarang anti-turunan dari $f$ pada $[a,b]$, maka
        \[\int_a^b f(x) \,dx = F(b) - F(a)\]
    \end{teorema}
\end{frame}

\end{document}
